\section*{الفصل \ref{ch:b2:affine} --- الفضاءات الأفينية}

\begin{solution}{exo:b2:affine:1}
$\{x + y + z = 1\}$: مستو أفيني، اتجاهه المستوي المتجهي $\{x
+ y + z = 0\}$، وبعده $2$. وبإضافة
$x - z = 3$: مستقيم أفيني (فالمعادلتان مستقلتان)، اتجاهه $\{x + y + z = 0,\ x = z\}
= \operatorname{Vect}\bigl((1, -2, 1)\bigr)$، وبعده $1$.
$\{x^2
+ y^2 = 1\}$: أسطوانة --- وهي غير مستقرة بمراكز الثقل (فمنتصف $(1,0,0)$ و$(-1,0,0)$
هو المبدأ، وهو خارج الأسطوانة): ومنه فهي ليست أفينية. وجملة متوافقة
$MX = B$: \omterm{def:b2:affine:def}{فضاء أفيني} جزئي $X_0 + \ker M$ بعده $\dim\ker M$، كما ذُكّر به في
\cref{def:b2:affine:subspace}.
\end{solution}

\begin{solution}{exo:b2:affine:2}
تعني $C = \operatorname{bar}(A, 1-t; B, t)$ أن $\vect{AC} =
t\,\vect{AB}$: ووجود $t$ من هذا النوع هو بالضبط ارتباط
$\vect{AC}$ مع $\vect{AB} \neq 0$، أي الاستقامة. والمواضع: $t
= \frac12$: المنتصف؛ $t = 2$: بعد
$B$، على مسافة منه تساوي مسافة $B$ ($\vect{AC} = 2\vect{AB}$)؛ $t = -1$: نظير $B$
بالتناظر المركزي حول $A$.
\end{solution}

\begin{solution}{exo:b2:affine:3}
$M$ هي مصفوفة الإسقاط على $\operatorname{Vect}(1,1)$ وفق $(1,-1)$ (تحقق من $M^2 = M$). صورة
$f$: $\{MX + C\} = C +
\operatorname{im} M$: أي المستقيم الأفيني المار بالنقطة $(1,0)$ والموجه
بالمتجهة $(1,1)$. النقط الثابتة: $X = MX + C$، أي $(I - M)X = C$؛ لكن $C =
(1, 0)^{\mathsf T}$ و$\operatorname{im}(I - M) =
\operatorname{Vect}(1,-1)$؛
فهل ينتمي $(1,0)$ إليه؟ إن $(1, 0) =
\alpha(1,-1)$ يفرض $\alpha = 1$ و$0 = -1$: كلا. فلا نقط
ثابتة. ثم
\[
f(f(X)) = M(MX + C) + C = MX + MC + C = f(X) + MC,
\qquad MC = \tfrac12(1,1)^{\mathsf T} :
\]
فالتطبيق $f\circ f$ هو $f$ متبوعا بانسحاب موازٍ لمستقيم الصورة ---
أي إن $f$ ``إسقاط منزلق'': إسقاط على المستقيم مركب مع انزلاق.
\end{solution}

\begin{solution}{exo:b2:affine:4}
$I = \operatorname{bar}(B, 2; C, 1)$ (لأن $\vect{BI} =
\frac13\vect{BC}$ يضع $I$ أقرب إلى $B$: بالوزن $2$ على
$B$، والوزن $1$ على $C$ --- وللتحقق: $\vect{BI} = \frac{1}{3}\vect{BC}$).
وبالمثل $J = \operatorname{bar}(C, 2; A, 1)$، $K =
\operatorname{bar}(A, 2; B, 1)$. وبجمع الجمل الموزونة الثلاث، يكون مركز ثقل
$(I, 1; J, 1; K, 1)$ (ووزن كل منها الكلي $3$، فنستبدل $I$ بجملته، وهكذا) هو
\[
\operatorname{bar}\bigl(A, 1 + 2;\ B, 2 + 1;\ C, 1 + 2\bigr)
= \operatorname{bar}(A, 1; B, 1; C, 1) = G ,
\]
أي مركز ثقل $ABC$: فللمثلث $IJK$ \omterm{def:b2:affine:barycenter}{مركز الثقل} نفسه --- وهو أمر متوقع،
لأن الإنشاء يعامل $A, B, C$ معاملة دائرية ولأن \omterm{def:b2:affine:barycenter}{مركز الثقل} هو النقطة
الثابتة الوحيدة للتناظر الدائري للأوزان.
\end{solution}

\begin{solution}{exo:b2:affine:5}
نضع $g(u) = f(O + u) - f(O)$ (بالعمل في $\R^n$ مشعَّعا عند $O$)، $g(0) = 0$.

\emph{الجمعية:} $\frac{(O + u) + (O + v)}{2} = O + \frac{u +
v}{2}$، ومنه يعطي الحفاظ على المنتصفات $g\bigl(\frac{u+v}{2}\bigr)
= \frac{g(u) + g(v)}{2}$؛ ومع
$v = 0$: $g(u/2) = g(u)/2$؛ وبالجمع بين الأمرين، $g(u + v) = 2g\bigl(\frac{u+v}{2}\bigr) = g(u) + g(v)$.

\emph{التجانس بالنسبة إلى $\Q$:} تعطي الجمعية $g(nu) = ng(u)$ ($n \in
\N$،
بالتراجع)، ثم $g(-u) = -g(u)$ (بالجمع)، ثم $g(\frac pq u) =
\frac pq g(u)$ (بتطبيق $q$ واستعمال تباين
الضرب).

\emph{التجانس بالنسبة إلى $\R$:} من أجل $t \in \R$، نأخذ أعدادا ناطقة
$t_n \to
t$: $g(t_nu) = t_ng(u)$، ويسمح اتصال $g$ (الموروث من $f$) بالمرور إلى
النهاية: $g(tu) = tg(u)$. ومنه فإن $g$ خطي و$f =
f(O) + g$: أي إنه أفيني.
\end{solution}

\begin{solution}{exo:b2:affine:6}
الجزء الخطي $\varphi(x,y) = (y, x)$: هو الانعكاس حول القطر $y = x$ (وهو متعامد ومحدده
$-1$). النقط الثابتة: يؤول $(x, y) = (y
+ 1, x - 1)$ إلى المعادلة الوحيدة $y = x - 1$
(فالمركبتان متكافئتان): ومنه فإن كل نقطة من المستقيم $y = x - 1$ ثابتة.
ومنه فإن $f$ يثبّت ذلك المستقيم نقطة نقطة: أي إن $f$ هو
\emph{الانعكاس} حول ذلك المحور (وهو تقايس له مستقيم من النقط الثابتة
وجزؤه الخطي انعكاس). وبانسجام مع ذلك، $f \circ
f(x,y) = f(y+1, x-1) = (x - 1 + 1, y + 1 - 1) = (x, y)$: أي إنه تقابل تقابلي
مع نفسه، كما يجب أن يكون كل انعكاس.
\end{solution}

\begin{solution}{exo:b2:affine:7}
المتجهات $\vect{A_1A_i}$ وعددها $n + 1$ ($i \geq 2$) مرتبطة في البعد $n$: فتوجد
$\mu_i$ ليست معدومة كلها تحقق $\sum_{i\geq2}
\mu_i\vect{A_1A_i} = 0$؛ ونضع $\mu_1 = -\sum_{i \geq 2}\mu_i$، ومنه $\sum_{i}\mu_i = 0$
و$\sum_i \mu_i\,\vect{OA_i} = 0$ من أجل كل $O$، مع كون $\mu_i$ ليست معدومة كلها. ونقسم
الأدلة: $P = \{i : \mu_i >
0\}$، $N = \{i : \mu_i < 0\}$، وكلتاهما غير خالية (فمجموع $\mu_i$ معدوم
وهي ليست معدومة كلها). وبوضع $s = \sum_{i\in P}\mu_i =
-\sum_{i \in N}\mu_i > 0$:
\[
\operatorname{bar}\bigl(A_i, \tfrac{\mu_i}{s}\bigr)_{i \in P}
= \operatorname{bar}\bigl(A_i, \tfrac{-\mu_i}{s}\bigr)_{i \in N},
\]
(فالطرفان يساويان النقطة $X$ مع $\vect{OX} = \frac1s\sum_{i\in
P}\mu_i\vect{OA_i}$، بحكم العلاقة): أي نقطة
مشتركة بين الغلافين المحدبين، بمجموعتي أدلة منفصلتين.
\end{solution}

\begin{solution}{exo:b2:affine:8}
\emph{الصورة تساوي مجموعة النقط الثابتة:} من أجل $Y = f(X)$، $f(Y) = f(f(X)) =
f(X) = Y$:
فكل نقطة صورة ثابتة؛ وبالعكس فالنقط الثابتة صور. ومنه فإن $\mathcal{F} = \operatorname{im} f =
\operatorname{Fix}(f)$ غير
خالية، وهي \omterm{def:b2:affine:def}{فضاء أفيني} جزئي (فهي صورة \omterm{def:b2:affine:subspace}{تطبيق أفيني})، اتجاهه $\operatorname{im}\vec f$.

\emph{بنية الإسقاط:} التطبيق $\vec f$ متساوي القوى ($\vec{f\circ
f} = \vec f^{\,2} = \vec f$)، ومنه
$E = \operatorname{im}\vec f
\oplus \ker \vec f$ (\cref{ex:b2:reduction:projections}). ومن أجل أي نقطة $X$، ننظر في المتجهة $\vect{f(X)\,X}$؛
وبتطبيق $\vec f$:
\[
\vec f\bigl(\vect{f(X)\,X}\bigr) = \vect{f(f(X))\,f(X)} = 0
\qquad (f \circ f = f),
\]
ومنه $\vect{f(X)\,X} \in \ker\vec f$. ومنه فإن $X = f(X) +
\vect{f(X)X}$ يُظهر $X$ نقطة من $\mathcal{F}$ منسحبة
بمتجهة من $\ker\vec f$: أي إن $f$ هو بالضبط الإسقاط على $\mathcal{F}$ وفق
$\ker\vec f$. وبالعكس فإن مثل هذه الإسقاطات تحقق $f \circ f = f$ بوضوح.
\end{solution}

\begin{solution}{exo:b2:affine:9}
لنضع $M = \operatorname{bar}(B, 2;\ C, 3)$، ووزنها الكلي $5$.
تعطي التجميعية $G = \operatorname{bar}(A, 1;\ M, 5)$، ومنه $\vect{AG} = \frac56\,\vect{AM}$: فالنقطة $G$ تقع على القطعة
$\intcc AM$ عند خمسة أسداسها انطلاقا من $A$. وبما أن $A \notin
(BC)$، فإن
المستقيم $(AG) = (AM)$ يقطع $(BC)$ في النقطة الوحيدة $M$، مع $\vect{BM} = \frac35\,\vect{BC}$.
وبالمثل، مع $N =
\operatorname{bar}(C, 3;\ A, 1)$ (بوزن كلي $4$، $\vect{CN}
= \frac14\,\vect{CA}$)، تعطي التجميعية
$G =
\operatorname{bar}(B, 2;\ N, 4)$: فالمستقيم $(BG)$ يقطع $(CA)$ في $N$، ولدينا $\vect{BG} = \frac46\,\vect{BN} =
\frac23\,\vect{BN}$.
\end{solution}

\begin{solution}{exo:b2:affine:10}
نتخذ $O$ مبدأ ونكتب النقط متجهات: $h_{\Omega, \lambda}(x) = \omega + \lambda(x - \omega)$ مع $\omega = \vect{O\Omega}$. والمركب
$g = h_{\Omega',
\mu} \circ h_{\Omega, \lambda}$ أفيني وجزؤه الخطي $\mu\lambda\,\mathrm{id}$. فإذا كان $\lambda\mu \neq 1$: فإن $1 \notin
\operatorname{Sp}(\lambda\mu\,\mathrm{id})$، ومنه
تعطي \cref{prop:b2:affine:fixedpoint} نقطة ثابتة وحيدة $\Omega''$، وبإعادة التمركز عندها، $g =
\lambda\mu\,\mathrm{id}$: أي
التحاكي $h_{\Omega'',
\lambda\mu}$. وإذا كان $\lambda\mu = 1$ فالجزء الخطي هو التطبيق المطابق،
ومنه فإن $g$ انسحاب؛ وبالنشر،
\[
g(x) = \omega' + \mu\bigl(\omega + \lambda(x - \omega) -
\omega'\bigr) = x + (1 - \mu)\,\omega' + \mu(1 -
\lambda)\,\omega .
\]
ومن أجل $\lambda = \mu = -1$ (تناظران مركزيان) تكون المتجهة $2\omega' - 2\omega = 2\,\vect{\Omega\Omega'}$: أي إن مركب
التناظر المركزي حول $\Omega$ ثم حول $\Omega'$ هو الانسحاب بالمتجهة $2\,\vect{\Omega\Omega'}$.
\end{solution}

\begin{solution}{exo:b2:affine:11}
تقول $\vect{A'B} = \alpha\,\vect{A'C}$ بالضبط إن $1\cdot
\vect{A'B} - \alpha\,\vect{A'C} = 0$، أي $A' =
\operatorname{bar}(B, 1;\ C, -\alpha)$ (بوزن كلي $1 -
\alpha \neq 0$ لأن
$B \neq C$)؛ وبالمثل $B' =
\operatorname{bar}(C, 1;\ A, -\beta)$ و$C' =
\operatorname{bar}(A, 1;\ B, -\gamma)$.

\emph{محك الاستقامة.} نعطي كل نقطة $P$ سطرها الثقلي الموحد $p = (p_A, p_B, p_C)$،
$p_A + p_B +
p_C = 1$، بالنسبة إلى $(A, B, C)$. فإذا كان $\sum_i c_i p_i = 0$ مع $(c_1, c_2, c_3) \neq 0$ من أجل ثلاث
نقط $P_1, P_2,
P_3$، فإن جمع المعاملات يعطي $\sum c_i = 0$، و$\sum_i c_i \vect{OP_i} = \sum_j \bigl(\sum_i
c_ip_{ij}\bigr)\vect{OV_j} = 0$: أي إن النقط
$P_i$ متعلقة \omterm{def:b2:affine:subspace}{أفينيا}، ومنه فهي على استقامة واحدة. وبالعكس، يعطي
تعلق أفيني $(t_i)$ المقدار $w = \sum t_ip_i$ بمعاملات مجموعها $0$ و$\sum_j w_j\vect{OV_j} = 0$؛
وبالنشر انطلاقا من $A$، $w_B
\vect{AB} + w_C\vect{AC} = 0$، ومنه $w = 0$ بحكم الاستقلال
الأفيني للنقط $(A, B, C)$: فالأسطر مرتبطة خطيا.
ومنه فإن الاستقامة تؤول إلى انعدام \omterm{def:b2:linalg:det}{محدد} $3 \times 3$، وضرب الأسطر في
العوامل غير المعدومة $1 - \alpha$ و$1 -
\beta$ و$1 - \gamma$ لا يغير شيئا:
\[
\det\begin{pmatrix} 0 & 1 & -\alpha\\ -\beta & 0 & 1\\ 1 &
-\gamma & 0\end{pmatrix}
= 1 - \alpha\beta\gamma .
\]
ومنه تكون $A', B', C'$ على استقامة واحدة إذا وفقط إذا كان $\alpha\beta\gamma = 1$: وهي
مبرهنة منيلاوس.
\end{solution}

\begin{solution}{exo:b2:affine:12}
لنضع $\Delta = \{\lambda \in \R^{n+1} : \lambda_i \geq 0,\
\sum\lambda_i = 1\}$: وهي مغلقة ومحدودة في $\R^{n+1}$، ومنه متراصة، والمجموعة
$K^{n+1}$ متراصة بوصفها جداء منتهيا. والتطبيق
\[
\Phi \colon \Delta \times K^{n+1} \to \R^n, \qquad
\Phi(\lambda, x_0, \dots, x_n) = \sum_{i=0}^n \lambda_i x_i
\]
\omterm{def:b2:metric:continuity}{متصل}، وتقول \cref{thm:b2:affine:caratheodory} بالضبط إن $\operatorname{conv}(K) = \Phi(\Delta \times
K^{n+1})$: أي صورة مجموعة متراصة بتطبيق
\omterm{def:b2:metric:continuity}{متصل} (\cref{thm:b2:metric:compactprops})، ومنه فهي متراصة.

أما في حالة مجموعة مغلقة: فلنأخذ $S = (\R \times \{0\}) \cup \{(0,
1)\}$، وهي مغلقة في $\R^2$.
والتوليفة المحدبة التي تضع الوزن $t$ على $(0,1)$ والوزن $1 - t$
على نقط المحور إحداثيتها الثانية $t$، ومنه
\[
\operatorname{conv}(S) = \bigl(\R \times \intco01\bigr) \cup
\{(0,1)\}
\]
(فمن أجل $0 \leq t < 1$، $(x, t) = t\,(0,1) + (1-t)\,\bigl(\tfrac
x{1-t}, 0\bigr)$). والنقطة $(1, 1) = \lim_{t \to 1}(1, t)$ ملتصقة بالغلاف لكنها ليست
فيه: فهو غير مغلق.
\end{solution}

\begin{solution}{pb:b2:affine:1}
\textbf{1.} نغير نقطة الأساس إلى $A_j$: من أجل $i \neq j$، $\vect{A_jA_i} =
\vect{A_0A_i} - \vect{A_0A_j}$.
فإذا كان $\sum_{i \neq j}
c_i\vect{A_jA_i} = 0$، أعطى النشر $\sum_{i \notin \{0,
j\}} c_i\,\vect{A_0A_i} - \bigl(\sum_{i\neq j}
c_i\bigr)\vect{A_0A_j} = 0$؛ ويفرض استقلال المقادير $\vect{A_0A_i}$
أن يكون $c_i = 0$ من أجل $i \notin \{0, j\}$، ثم $c_0 = 0$: أي الاستقلال عند $A_j$.
أما التكافؤ: فعائلتا أوزان $(\lambda_i)$ و$(\lambda_i')$ مجموع كل منهما $1$
ولهما \omterm{def:b2:affine:barycenter}{مركز الثقل} نفسه تعطيان، مع $\nu = \lambda -
\lambda'$: $\sum\nu_i = 0$ و(بالمبدأ $A_0$)
$\sum_{i \geq
1}\nu_i\,\vect{A_0A_i} = 0$، ومنه $\nu = 0$ بحكم الاستقلال.
وبالعكس، فإن علاقة غير تافهة $\sum_{i\geq1}\mu_i
\vect{A_0A_i} = 0$، مكملة بالمقدار $\mu_0 = -\sum_{i\geq1}
\mu_i$، تسمح
بإضافة $t(\mu_i)$ إلى أي عائلة أوزان دون تحريك \omterm{def:b2:affine:barycenter}{مركز الثقل}: أي انعدام
الوحدانية.

\textbf{2.} $(\vect{AB}, \vect{AC})$ أساس للفضاء $\R^2$:
نكتب $\vect{AM} = \beta\,\vect{AB} + \gamma\,\vect{AC}$ (بصورة وحيدة) ونضع $\alpha = 1 - \beta - \gamma$؛ وشرط \omterm{def:b2:affine:barycenter}{مركز الثقل} بالمبدأ
$A$ يُكتب بالضبط $\vect{AM} =
\beta\,\vect{AB} + \gamma\,\vect{AC}$. والوحدانية هي السؤال 1.

\textbf{3.} انطلاقا من $\alpha\vect{MA} + \beta\vect{MB} +
\gamma\vect{MC} = 0$ وعلاقة شال، $\vect{MA} =
-\beta\,\vect{AB} - \gamma\,\vect{AC}$. ومع $D =
\det(\vect{AB}, \vect{AC})$:
\begin{align*}
\det(\vect{MB}, \vect{MC})
&= \det(\vect{MA} + \vect{AB},\ \vect{MA} + \vect{AC})\\
&= \det(\vect{MA}, \vect{AC}) + \det(\vect{AB}, \vect{MA}) +
D = -\beta D - \gamma D + D = \alpha D,
\end{align*}
باستعمال ثنائية الخطية و$\det(\vect{MA}, \vect{AC}) = -\beta D$، $\det(\vect{AB}, \vect{MA}) = -\gamma D$. وتنتج الصيغتان الأخريان
بالحساب نفسه مع تبديل الأدوار دائريا.

\textbf{4.} بالتعريف، $\operatorname{conv}\{A, B, C\}$ هي مجموعة مراكز الثقل ذات الأوزان
الموجبة؛ وبتوحيد الأوزان لتصير $1$ واستدعاء الوحدانية (السؤال 2)،
يكون $M \in \operatorname{conv}\{A,B,C\}$ إذا وفقط إذا كان $\alpha, \beta,
\gamma \geq 0$. وكل إحداثية دالة أفينية بدلالة
$M$ (السؤال 3: \omterm{def:b2:linalg:det}{محدد} من الحجم $2\times2$ أحد أعمدته أفيني بدلالة
$M$)، ومنه فإن كل شرط إشارة مفتوح يعرّف نصف مستو مفتوح. والنمط
$(-,-,-)$ يناقض $\alpha + \beta
+ \gamma = 1$؛ أما كل من الأنماط السبعة الباقية فمحقق: يكفي
ضرب ثلاثية تحترم الإشارات وفيها معامل $+$ واحد على الأقل بحيث
يصير المجموع (الموجب) $1$ --- مثلا $(-1, 1, 1)$، $(3, -1, -1)$، $(\frac13, \frac13, \frac13)$،
وتبديلاتها.

\textbf{5.} يحافظ \omterm{def:b2:affine:subspace}{التطبيق الأفيني} على مراكز الثقل (\cref{def:b2:affine:subspace})، ومنه
$u(M) = \alpha u(A) +
\beta u(B) + \gamma u(C)$. وبكتابة $u(x, y) = ax + by + c$ مع $(a, b) \neq (0,0)$: تكون $\{u = c'\}$ مستقيما، وكل مستقيم
$ax + by = c'$ مجموعة مستوى من هذا النوع. وإذا كان $u(M), u(N) \geq
c$ و$t \in \intcc01$، فإن
$u\bigl(\operatorname{bar}(M,
1-t; N, t)\bigr) = (1-t)u(M) + tu(N) \geq c$: ومنه فإن أنصاف المستويات محدبة.

\textbf{6.} تعطي الشروط $\sum\mu_i = 0$ و$3\mu_2 + \mu_4 =
0$ و$3\mu_3 + \mu_4 = 0$ (بأخذ $\mu_4 = 3$)
التعلق $(\mu_1, \mu_2, \mu_3, \mu_4) = (-1, -1, -1, 3)$.
وتنقسم الإشارات على الصورة $\{A_4\} \mid \{A_1, A_2, A_3\}$، وبتوحيد كل طرف بالمقدار $3$:
\[
A_4 = \operatorname{bar}\bigl(A_1, \tfrac13;\ A_2, \tfrac13;\
A_3, \tfrac13\bigr) = (1,1),
\]
أي مركز ثقل المثلث: فنقطة رادون هي $A_4$ نفسها، وهي تقع فعلا داخل
المثلث $A_1A_2A_3$.

\textbf{7.} التطبيق الخطي $\Phi \colon \R^{n+2} \to \R
\times \R^n$، $\mu \mapsto (\sum\mu_i,\
\sum\mu_i\vect{OA_i})$، رتبته $\leq n + 1$، ومنه
$\dim\ker
\Phi \geq 1$. فلو وُجد تعلقان مستقلان $\mu, \mu'$، لكانت توليفة مناسبة $\nu = \mu'_{n+2}\mu -
\mu_{n+2}\mu'$
(أو $\mu$ نفسه إذا انعدم المعاملان الأخيران معا) تعلقا \emph{غير
معدوم} يحقق $\nu_{n+2}
= 0$؛ وبالقصر على $A_1, \dots, A_{n+1}$ وتغيير الأساس إلى $A_1$،
يوجد $\nu_i \neq 0$ مع $i \geq 2$ (فوزن واحد غير معدوم لا يمكن أن يكون مجموعه
معدوما)، وهو ما يعطي علاقة غير تافهة $\sum_{i\geq2}\nu_i\vect{A_1A_i} = 0$: فتكون النقط $n+1$
متعلقة \omterm{def:b2:affine:subspace}{أفينيا}، خلافا للوضع العام. ومنه $\dim\ker
\Phi = 1$. وتبين حجة القصر
نفسها أن التعلق غير المعدوم لا يحوي أي معامل معدوم.

\textbf{8.} ليكن $\mu \neq 0$ تعلقا، ولنضع $P = \{i :
\mu_i > 0\}$ و$N = \{i : \mu_i < 0\}$: وكلتاهما غير
خالية ($\sum\mu_i = 0$، $\mu \neq 0$) وشاملة (فلا معامل معدوم). وينتج إنشاء رادون
(\cref{exo:b2:affine:7}) نقطة الغلاف المشتركة انطلاقا من هذا التقسيم بالضبط.
وبما أن التعلق وحيد إلى غاية عدد غير معدوم (السؤال 7)، فإن الزوج
غير المرتب $\{P, N\}$ --- ومنه تقسيم رادون --- وحيد.

\textbf{9.} الكتلتان غير خاليتين، ومنه فالنمط $(1,3)$ أو $(2,2)$.
النمط $(1,3)$، بالكتلة $\{j\}$: تقع نقطة رادون في $\operatorname{conv}\{A_j\} = \{A_j\}$، ومنه فهي
$A_j \in
\operatorname{conv}$ للنقط الثلاث الأخرى؛ ولا يمكن أن تقع على ضلع (وإلا كانت ثلاث
نقط على استقامة واحدة، خلافا للوضع العام)، ومنه فإن $A_j$ داخلية
للمثلث. والنمط $(2,2)$، بالكتلتين $\{i,j\} \mid \{k,l\}$: تقع نقطة رادون $z$ على
$\intcc{A_i}{A_j} \cap \intcc{A_k}{A_l}$، وليست $z$ طرفا (وإلا كانت ثلاث نقط على استقامة واحدة):
ومنه تتقاطع القطعتان في نقطة داخلية. وعلاوة على ذلك لا تقع أي نقطة
في غلاف الأخريات: فاحتواء من هذا النوع $A_l =
\operatorname{bar}(A_i, \lambda_i)_{i \neq l}$ مع $\lambda_i \geq 0$ تعلق أفيني
نمط إشاراته $(+,+,+,-)$، وهو ما كان سيجعل التقسيم $(1,3)$ بحكم الوحدانية
(السؤال 8). ومنه ففي الحالة $(2,2)$ تكون النقط الأربع في وضع محدب
والقطعتان المتقاطعتان هما القطران.

\textbf{10.} تعطي المعادلات $\mu_2 + \mu_3 = 0$ و$\mu_3 +
\mu_4 = 0$ و$\sum\mu_i = 0$ التعلق $(1, -1, 1,
-1)$:
بالتقسيم $\{(0,0), (1,1)\} \mid \{(1,0), (0,1)\}$، و
\[
\operatorname{bar}\bigl((0,0), \tfrac12;\ (1,1),
\tfrac12\bigr) = \bigl(\tfrac12, \tfrac12\bigr) =
\operatorname{bar}\bigl((1,0), \tfrac12;\ (0,1),
\tfrac12\bigr) :
\]
فنقطة رادون هي مركز المربع، حيث يتقاطع القطران --- أي النمط $(2,2)$،
كما تتنبأ الصورة.

\textbf{11.} تعطي مبرهنة رادون مطبقة على $x_1, \dots, x_4$ الكتلتين $I \mid J$
ونقطة $z \in \operatorname{conv}\{x_i : i
\in I\} \cap \operatorname{conv}\{x_j : j \in J\}$. نثبت $k
\in \{1, \dots, 4\}$، ولتكن $k \in I$. وكل $j \in J$ يحقق $j \neq k$،
ومنه $x_j \in C_k$ بحكم الاختيار $x_j \in
\bigcap_{l \neq j}C_l$؛ وبما أن $C_k$ محدبة، فإن $z \in
\operatorname{conv}\{x_j : j \in J\} \subseteq C_k$.
وبما أن $k$ كيفما كانت، فإن $z \in C_1 \cap C_2 \cap C_3 \cap C_4$.

\textbf{12.} بالتراجع على $m$. من أجل $m = 3$ تكون الفرضية هي
النتيجة؛ والحالة $m = 4$ هي السؤال 11. ليكن $m \geq 4$، ولنفترض النص من
أجل $m$ مجموعة، ولنأخذ $C_1, \dots,
C_{m+1}$ بخاصية التقاطع الثلاثي. نضع $C_m' =
C_m \cap C_{m+1}$
وهي محدبة. وللعائلة $C_1, \dots, C_{m-1},
C_m'$ عدد من العناصر يساوي $m$؛ وكل ثلاثية
تتجنب $C_m'$ تتقاطع بحكم الفرضية، وكل ثلاثية $\{C_i, C_j, C_m'\}$ تقاطعها $C_i \cap C_j \cap C_m \cap C_{m+1}$،
وهو غير خال بحسب السؤال 11 مطبقا على $C_i, C_j, C_m, C_{m+1}$ (فكل ثلاث منها تتلاقى،
بحكم الفرضية). وتعطي فرضية التراجع الآن نقطة مشتركة للعائلة الجديدة،
أي لجميع المجموعات $m+1$.

\textbf{13.} (1) الأضلاع المغلقة $\intcc PQ$ و$\intcc QR$ و$\intcc RP$ لمثلث غير
منحل: يتقاسم كل ضلعين رأسا، لكن نقطة مشتركة بين الثلاثة كانت ستقع
في $\intcc
PQ \cap \intcc RP = \{P\}$ وفي $\intcc QR$، وهو ما يستبعد $P$. (2) كل ثلاث من
المجموعات $S_i = \{x_1, \dots, x_4\}
\setminus \{x_i\}$ تُسقط ثلاثا من النقط الأربع، فتبقي نقطة مشتركة
واحدة بالضبط؛ أما التقاطع الكلي فيسقط كل نقطة. والمجموعات $S_i$
منتهية لا محدبة: فالتحدب جوهري. (3) عدد منته من المجموعات $H_k = \intco{k}{+\infty} \times
\R$
يتقاطع في $\intco{k_{\max}}{+\infty} \times \R \neq
\emptyset$، ومع ذلك لا توجد نقطة تحقق $x \geq k$ من أجل كل $k \in \N$:
ففي حالة العائلات اللانهائية، يكون التراص جوهريا.

\textbf{14.} لنفترض $\bigcap_{i \in I}K_i = \emptyset$ ولنثبّت $i_0$. كل $x \in K_{i_0}$ يفوّت مجموعة
$K_i$ ما، ومنه $K_{i_0} \subseteq \bigcup_{i \in I}(\R^2 \setminus K_i)$، وهي تغطية بمفتوحات (فالمجموعة $K_i$ متراصة،
ومنه مغلقة). وبخاصية بوريل--لوبيغ (\cref{thm:b2:metric:borellebesgue}) يكفي عدد منته منها:
$K_{i_0} \cap K_{i_1} \cap \dots \cap K_{i_N} =
\emptyset$. لكن كل ثلاثة عناصر من هذه العائلة المنتهية من \omterm{def:b2:affine:convex}{المجموعات
المحدبة} تتقاطع، ومنه يجعل السؤال 12 التقاطع غير خال: وهو تناقض.

\textbf{15.} نضع $D_p = \overline D(p, r)$ من أجل $p \in S$: وهي مجموعات محدبة متراصة.
ومن أجل $p, q, s \in S$، تعطي الفرضية قرصا مغلقا $\overline D(z, r)$ يحوي $p, q, s$؛ ومنه
$\norm{\vect{zp}}, \norm{\vect{zq}}, \norm{\vect{zs}}
\leq r$، أي $z \in D_p \cap D_q \cap D_s$. وبمبرهنة هيلي (السؤال 12؛ فالعائلة منتهية) توجد
$c \in
\bigcap_{p \in S}D_p$: فكل $p \in S$ يحقق $\norm{\vect{cp}} \leq r$، ومنه $S \subseteq \overline D(c,
r)$.

\textbf{16.} لنضع $c = x_{\lceil n/2\rceil}$. ونصف المستقيم المغلق الحاوي $c$ هو
$\intoc{-\infty}{t}$ مع $t \geq
c$ أو $\intco{t}{+\infty}$ مع $t \leq c$. والأول يحوي $x_1, \dots, x_{\lceil n/2\rceil}$: أي ما
لا يقل عن $\lceil n/2\rceil \geq n/2$ نقطة. والثاني يحوي $x_{\lceil n/2\rceil}, \dots, x_n$: أي $n - \lceil
n/2\rceil + 1 = \floor{n/2} + 1 > n/2$ بالضبط.

\textbf{17.} $\abs{A \cap B} = \abs A + \abs B - \abs{A \cup
B} \geq \abs A + \abs B - n > \tfrac{4n}3 - n = \tfrac n3$، ثم
\[
\abs{A \cap B \cap C} \geq \abs{A \cap B} + \abs C - n >
\tfrac n3 + \tfrac{2n}3 - n = 0 .
\]

\textbf{18.} لاحظ أن $m > 2n/3$. ومن أجل $T_1, T_2, T_3 \subseteq S$ من الحجم $m$، يوفر
السؤال 17 نقطة $x \in T_1
\cap T_2 \cap T_3$؛ ومنه $x \in \operatorname{conv}(T_i)$ من أجل كل $i$: أي إن كل
ثلاثة عناصر من $\mathcal F$ تتلاقى. والعائلة منتهية (فعدد المجموعات
الجزئية من $S$ منته) وتتكون من مجموعات محدبة، ومنه تعطي مبرهنة
هيلي (السؤال 12) المقدار $c \in
\bigcap_{\abs T = m}\operatorname{conv}(T)$.

\textbf{19.} لنفترض أن نصف مستو مغلق $H \ni c$ يحوي أقل من $n/3$ نقطة
من $S$. متممه $U$ نصف مستو مفتوح ومحدب، مع $\abs{S \cap U} > 2n/3$، ومنه
$\abs{S \cap U} \geq \floor{2n/3} + 1 = m$؛ نختار $T
\subseteq S \cap U$ بحيث $\abs T = m$. عندئذ $\operatorname{conv}(T) \subseteq U$ بحكم تحدب $U$،
ومنه $c \in \operatorname{conv}(T) \subseteq U$: وهو تناقض مع $c \in H$. ومنه فإن كل نصف مستو مغلق يحوي
$c$ يحوي ما لا يقل عن $n/3$ نقطة: أي إن $c$ نقطة مركز.

\textbf{20.} نأخذ المثلث متساوي الأضلاع ضلعه $L$ و$\varepsilon = L/100$. ولتكن
$c$ نقطة كيفما كانت؛ نُظهر نصف مستو مغلق يحوي $c$ ولا يحوي
أكثر من $k$ نقطة.

\emph{الحالة 1: تقع $c$ على مسافة لا تتجاوز $L/10$ من رأس، وليكن
$B$.} الاتجاهان من $c$ نحو $A$ ونحو $C$ ينحرفان عن
الاتجاهين $B \to A$ و$B \to C$ بما لا يتجاوز $\arcsin\bigl(\tfrac{L/10}{9L/10}\bigr) = \arcsin\tfrac19$، ومنه فهما يصنعان
زاوية $\leq \tfrac\pi3 + 2\arcsin\tfrac19 <
\tfrac{2\pi}3$. \emph{الحالة 2: تقع $c$ على مسافة $> L/10$ من جميع
الرؤوس.} فإذا كانت $c$ داخل المثلث، فإن الفجوات الزاوية الثلاث
بين الاتجاهات $u_A, u_B, u_C$ من $c$ نحو الرؤوس مجموعها $2\pi$، ومنه
فإحدى الفجوات $\leq 2\pi/3$؛ وإذا كانت $c$ خارجه، فإن الاتجاهات الثلاثة
تقع في نصف مستو مفتوح من الاتجاهات ويصنع اثنان منها زاوية $<
\pi/2$.
وفي جميع الأحوال يصنع اتجاهان، وليكونا نحو $X$ و$Y$، زاوية
$\leq 2\pi/3$؛ وليكن $w$ منصفهما الواحدي، بحيث $\langle u_X, w\rangle, \langle u_Y, w\rangle \geq
\cos\tfrac\pi3 = \tfrac12$. ومن أجل أي نقطة
$b$ من القرص المحيط بالنقطة $X$:
\[
\langle \vect{cb}, w\rangle \geq
\tfrac12\norm{\vect{cX}} - \varepsilon > 0,
\]
لأن $\norm{\vect{cX}} \geq L/10 > 2\varepsilon$ (وبالمثل من أجل $Y$): ومنه فإن نصف المستوي المفتوح
$\{\langle \vect{cx},
w\rangle > 0\}$ يبتلع العنقودين معا. ومتممه المغلق يحوي $c$ ولا يحوي أكثر
من نقط العنقود الثالث وعددها $k$. ومنه فلا تتفوق أي نقطة من
المستوي على $n/3$: ومع السؤال 19، يكون ثابت نقطة المركز $1/3$
بالضبط.

\textbf{21.} نرتب الزوايا؛ فالكبرى، $\theta$، تحقق $\theta \geq \pi/3$ (فمجموع
الثلاث $\pi$). فإذا كان $\theta \geq \pi/2$، وليكن في $R$، فلتكن $M$
منتصف الضلع المقابل $\intcc PQ$. وتعطي صيغة المتوسط ($\vect{RM}
= \tfrac12(\vect{RP} + \vect{RQ})$، بالنشر
وحذف $\langle\vect{RP}, \vect{RQ}\rangle$ بقانون جيوب التمام)
\[
\norm{\vect{RM}}^2 = \tfrac12\norm{\vect{RP}}^2 +
\tfrac12\norm{\vect{RQ}}^2 - \tfrac14\norm{\vect{PQ}}^2 \leq
\tfrac12\norm{\vect{PQ}}^2 - \tfrac14\norm{\vect{PQ}}^2 =
\tfrac14\norm{\vect{PQ}}^2,
\]
باستعمال $\norm{\vect{PQ}}^2 = \norm{\vect{RP}}^2 +
\norm{\vect{RQ}}^2 - 2\langle\vect{RP}, \vect{RQ}\rangle \geq
\norm{\vect{RP}}^2 + \norm{\vect{RQ}}^2$ (فالجداء السلمي $\leq 0$). ومنه فإن القرص الذي قطره
$\intcc PQ$، ونصف قطره $\leq \tfrac12 < \tfrac1{\sqrt3}$، يحوي النقط الثلاث (والثلاثيات المنحلة
المستقيمة تقع تحت $\theta = \pi$). وإذا كان $\theta < \pi/2$ فالمثلث حاد الزوايا؛
وبقانون الجيوب يكون نصف قطر الدائرة المحيطة $R_c = a/(2\sin\theta)$ مع $a \leq 1$ الضلع
المقابل للزاوية $\theta$، ويعطي $\theta \in
\intco{\pi/3}{\pi/2}$ المقدار $\sin\theta \geq \sqrt3/2$، ومنه $R_c \leq 1/\sqrt3$:
فالقرص المحيط يفي بالغرض.

\textbf{22.} من أجل $p \in S$ نضع $K_p = \overline D(p,
1/\sqrt3)$: وهي مجموعات محدبة متراصة.
وكل ثلاث نقط $p, q, s$ من $S$ المسافات بينها مثنى مثنى $\leq 1$،
ومنه يعطي السؤال 21 قرصا نصف قطره $1/\sqrt3$ يحويها: ومركزه يقع في
$K_p \cap K_q \cap K_s$. وبمبرهنة هيلي المتراصة (السؤال 14، فالعائلات الكيفية مسموح
بها) توجد $c \in \bigcap_{p\in
S}K_p$: فكل $p \in S$ يقع على مسافة لا تتجاوز $1/\sqrt3$ من
$c$، أي $S
\subseteq \overline D(c, 1/\sqrt3)$. وهذه هي مبرهنة يونغ في المستوي.

\textbf{23.} مع $G$ \omterm{def:b2:affine:barycenter}{مركز الثقل}، $\sum_i\vect{GA_i} = 0$، ومنه
\[
\sum_i\norm{\vect{OA_i}}^2 = \sum_i\norm{\vect{OG} +
\vect{GA_i}}^2 = 3\norm{\vect{OG}}^2 + 2\Bigl\langle
\vect{OG}, \sum_i\vect{GA_i}\Bigr\rangle +
\sum_i\norm{\vect{GA_i}}^2,
\]
وينعدم الحد الأوسط: وهي متطابقة ليبنتز. ومن أجل المثلث المتساوي
الأضلاع ذي الضلع $1$، لدينا $\norm{\vect{GA_i}} =
1/\sqrt3$ (أي ثلثا الارتفاع $\sqrt3/2$)،
ومنه $\sum_i\norm{\vect{GA_i}}^2 = 1$. فإذا حوى $\overline D(O, r)$ الرؤوس، فإن $3r^2 \geq
\sum_i\norm{\vect{OA_i}}^2 = 3\norm{\vect{OG}}^2 + 1 \geq 1$: أي $r \geq 1/\sqrt3$، ويفرض
التساوي أن يكون $O = G$ وأن تتساوى المسافات الثلاث بالمقدار $r$
--- وهو القرص المحيط. ومنه فإن ثابت يونغ $1/\sqrt3$ أمثلي.

\textbf{24.} \emph{مبرهنة هيلي في $\R^n$: إذا كانت $C_1, \dots, C_m$ ($m \geq n + 1$)
مجموعات جزئية محدبة من $\R^n$ وكانت كل $n + 1$ منها تتقاطع، فإنها
تتقاطع كلها.} الحالة الأساسية $m = n +
2$: نأخذ $x_i \in \bigcap_{j\neq i}C_j$؛ وتقسم مبرهنة
رادون (\cref{exo:b2:affine:7}) المجموعة $x_1, \dots, x_{n+2}$ إلى كتلتين $I \mid J$ لهما نقطة غلاف
مشتركة $z$، ومن أجل كل $k$ تتكون الكتلة التي لا تحوي $k$
من نقط من $C_k$، ومنه $z \in C_k$ بحكم التحدب، تماما كما في السؤال 11.
وخطوة التراجع من أجل $m \geq n + 2$: نستبدل $C_m,
C_{m+1}$ بالمجموعة $C_m \cap C_{m+1}$؛
وثلاثية من العائلة الجديدة تحوي العنصر المتقاطع تؤول إلى $(n+1)$
منظومة من $n + 2$ من المجموعات القديمة، وتعالجها الحالة الأساسية،
أما الثلاثيات الأخرى فتغطيها الفرضية. ونختم بفرضية التراجع.

\textbf{25.} يدخل \omterm{def:b2:structures:algebra}{الجبر} الخطي مرة واحدة: فالمتجهات وعددها $n + 2$ في
الفضاء ذي البعد $(n+1)$ من الأزواج (الوزن الكلي، الموضع الموزون) يجب
أن تكون مرتبطة --- وذلك هو التعلق الأفيني. وإشارات معاملاته تقسم
النقط إلى كتلتي رادون وتحول علاقة خطية واحدة إلى مساواة بين مركزي
ثقل موجبي الأوزان. ويُستعمل التحدب مرتين بالضبط: في خطوة هيلي (فغلاف
نقط $C_k$ يبقى في $C_k$) وفي التطبيقات (فأنصاف المستويات والأقراص
محدبة). ولم يدخل بعد المستوي إلا عبر عدد النقط $4 = 2 + 2$ المدفوعة إلى
مبرهنة رادون، أي عبر ``$3 = 2 + 1$'' في فرضية هيلي؛ وكل ما عدا ذلك كان
مستقلا عن البعد، كما يؤكد السؤال 24. وفي $\R^n$ تصير الثوابت: عدد
هيلي $n + 1$؛ وثابت نقطة المركز $\frac1{n+1}$ (فلكل مجموعة منتهية نقطة
يحوي كل نصف فضاء مغلق يمر بها جزءا $\geq
\frac1{n+1}$ منها)؛ ونصف قطر يونغ
$\sqrt{\frac{n}{2(n+1)}}$ من أجل المجموعات ذات القطر $1$ --- وهو يساوي $1/\sqrt3$
عندما $n = 2$.
\end{solution}
